\chapter{Development and Engineering Challenges}
\label{ch:development}

\section{From prototype to simulator}

Development began in early April 2026 with the smallest experiment that could fail usefully: a
cylinder, one oversized engine and a policy asked to reach the ground slowly. By the middle of
April it could perform a basic landing. That result justified continuing, but every simplifying
assumption was still visible. The engine had no realistic minimum throttle, vehicle mass never
changed, and one \texttt{FalconAgent} script owned nearly the entire simulation.

The next step was not one large jump to realism. I replaced assumptions one at a time and let
each addition reveal the next missing layer. Moving to Falcon-class dimensions led to multiple
engines, grid fins and RCS. Grid fins were meaningless without airflow, so drag and atmosphere
followed. Wind then required air-relative velocity, while fuel consumption required changing
mass and inertia. Finally, spool, slew and engine timing prevented the policy from treating the
actuators as instantaneous mathematical commands.

Those physical additions also changed the policy interface. Manually maintaining ML-Agents
observation and action sizes became error-prone, so a launch manager began deriving the schema
and creating parallel areas. That exposed a startup-order bug: ML-Agents connected before valid
runtime rockets existed. Solving it produced the dummy/actual-prefab sequence described in
Chapter~\ref{ch:architecture}.

The largest transition came with the right-side configuration panel. Parameters that had lived
beside the code using them now had to travel from an editable UI into a running experiment. This
forced the snapshot architecture. Adding hover, tracking and landing then forced a second split,
because spawn distributions, rewards, curricula, terminations and evaluation criteria could no
longer remain global switches inside one agent. By this point, the project had changed from a
landing demonstration into the configurable experiment system described in the previous three
chapters.

\section{Challenge 1: reward loopholes}

Reward development produced the clearest examples of why the simulator had to become observable.
A coefficient could look reasonable during manual testing and remain apparently successful for
millions of steps, until the policy discovered a repeated action sequence with a better return.
The most important behaviours were:

\begin{itemize}
  \item hovering near the ground instead of landing because proximity rewards continued to
  accrue;
  \item touching down with only one foot while satisfying an overly broad contact condition;
  \item switching off thrust and dropping when future return appeared worse than termination;
  \item remaining nearly inactive when movement costs outweighed progress; and
  \item flipping or forcing a collision to end an unfavourable episode early.
\end{itemize}

At first I responded by changing coefficients or adding special cases. That worked for one
behaviour and usually created another. The more useful interpretation was that each policy had
found a counterexample to the written reward specification. The long-term response was to expose
rates, events, terminal values, scales and thresholds in the scenario configuration, log their
individual contributions, and make termination order explicit. This does not prevent a poor
reward, but it makes the reason for a high return inspectable and allows controlled changes
without rewriting task code.

\section{Leg-landing problem-solving campaign}
\label{sec:landing-iterations}

The numbered L1--L11 series turned that lesson into a repeatable process. I started a new
iteration only after inspecting the previous policy in the simulator and in episode- and
step-level telemetry. Each run therefore begins with a concrete question left by the one before
it. Because duration, reward scale, actuator timing and curriculum difficulty changed, their raw
returns are not a controlled algorithm benchmark. Their value lies in the chain of evidence:
what failed, how the cause was isolated, and whether the next intervention removed it.

\begin{figure}[htbp]
  \centering
  \begin{tikzpicture}[
      node distance=0.35cm,
      phase/.style={draw=rocketNavy!35,rounded corners=3pt,align=center,text width=0.198\textwidth,
        minimum height=2.2cm,inner sep=5pt,font=\small,text=rocketInk},
      arrow/.style={-{Latex[length=2.2mm]},thick,draw=rocketBlue}]
    \node[phase,fill=rocketRed!7] (p1) {\textbf{L1--L4}\\Infrastructure, spawn physics and termination\\
      \emph{Make a physical attempt preferable to a cheap reset}};
    \node[phase,fill=rocketAmber!10,right=of p1] (p2) {\textbf{L5--L7}\\Contact quality and reward economics\\
      \emph{Make a soft, centred touchdown preferable to a crash}};
    \node[phase,fill=rocketTeal!9,right=of p2] (p3) {\textbf{L8--L9}\\Learnability and actuator timing\\
      \emph{Bootstrap landing, then restore constraints}};
    \node[phase,fill=rocketBlue!8,right=of p3] (p4) {\textbf{L10--L11}\\Efficiency and generalization\\
      \emph{Reward the outcome and test the full range}};
    \draw[arrow] (p1) -- (p2);
    \draw[arrow] (p2) -- (p3);
    \draw[arrow] (p3) -- (p4);
  \end{tikzpicture}
  \caption{The four phases of the L1--L11 landing campaign. Each phase begins with a policy or
  infrastructure failure and ends with a narrower unsolved problem.}
  \label{fig:landing-phases}
\end{figure}

\footnotesize
\begin{longtable}{p{0.055\textwidth}p{0.245\textwidth}p{0.27\textwidth}p{0.30\textwidth}}
  \caption{Landing campaign at a glance. The table gives the decision made after each run; the
  evidence and comparisons are developed in the phase narratives below.}
  \label{tab:landing-iterations}\\
  \toprule
  Run & Failure exposed & Decisive change & Result or next question \\
  \midrule
  \endfirsthead
  \toprule
  Run & Failure exposed & Decisive change & Result or next question \\
  \midrule
  \endhead
  L1 & Trainer never connected. & Clear the occupied Unity worker. & Infrastructure failure;
  no policy was tested. \\
  L2 & The rocket spawned without its prescribed downward velocity. & Apply velocity only after
  the rigid body becomes dynamic. & Escape disappeared as an easy interpretation of the task. \\
  L3 & The policy reached the pad, then crashed at about \SI{40}{m/s}. & Remove time pressure and
  extend impact severity beyond its low cap. & Touchdown became preferable to a fast reset. \\
  L4 & A vehicle could approach the pad and then leave after an off-pad collision. & Lock
  propulsion and classify the first external collision immediately. & The surface could no
  longer be used as a powered recovery manoeuvre. \\
  L5 & Almost every episode touched down, but most were hard or waited off-centre. & Add settled
  off-centre termination and stronger stability shaping. & The next problem became reward
  ordering, not contact frequency. \\
  L6 & A hard crash was cheaper than a careful miss. & Rebalance terminal outcomes and expose
  additive readiness components. & The economics improved, but navigation progress still leaked. \\
  L7 & Descending increased ``readiness'' even while lateral error grew. & Remove the altitude
  gate derivative from navigation progress. & The reward finally measured motion toward the
  centre. \\
  L8 & The full actuator delay made useful landing behaviour difficult to discover. & Temporarily
  ease startup, minimum-run and cooldown timing. & Robust landing emerged, together with roughly
  30 ignitions per success. \\
  L9 & Rapid engine cycling synthesized thrust below the minimum. & Restore part of the timing
  constraint and warm-start from L8. & Cycling fell, but efficiency and engine choice remained
  weak. \\
  L10 & Directly prescribing an engine would overfit the solution. & Start fresh and reward
  efficiency only after successful touchdown. & 151/250 fixed-band successes; full difficulty
  and cycling remained unsolved. \\
  L11 & One-engine continuous descent was still physically awkward. & Lower minimum throttle and
  strengthen success-only cycling cost. & 212/250 successes, including 18/50 at full difficulty;
  mean restarts fell to 1.56. \\
  \bottomrule
\end{longtable}
\normalsize

\subsection{L1--L4: make the episode physically meaningful}

The first phase prevented the policy from optimizing simulator and termination defects. L1
never began because the communication worker was occupied. L2 then appeared to learn quickly:
mean return improved from about $-18.7$ to $-17.2$, the vehicle centred to roughly
\SI{0.64}{m}, ignited all three engines and climbed out of the episode in about \SI{3.6}{s}.
This was not landing progress. The policy had found the cheapest reset, and the simulator had
silently removed the prescribed 10--\SI{18}{m/s} initial downward velocity.

After staged velocity application in L3, the policy could no longer begin from rest. It learned
to navigate near the pad and remain upright, but then converted the next reward defect into a
fast vertical crash. The hard-touchdown severity term had a low cap, while every additional
second still incurred cost. Removing that time pressure and extending severity in L4 made
physical contact more attractive. Inspection then revealed a different safety gap: vehicles
that lost control after approaching the pad could continue thrusting after an off-pad impact.
The response was a physical propulsion lockout and immediate classification on every first
external collision.

\subsection{L5--L7: make safe touchdown cheaper than failure}

Once almost every trajectory reached the surface, the target moved from contact quantity to
contact quality. L5 exposed two post-contact loopholes. A rocket could rock between feet for
several seconds before becoming stable, or remain calm just outside the success radius until a
discounted timeout. The settled off-centre termination added for L6 removed the second behaviour,
but L6 regressed to 1.6\% late-window success because the reward ordering was wrong: a fast hard
contact was much cheaper than a carefully controlled miss.

L7 corrected the ordering and made readiness additive rather than a product that became zero if
one component was poor. At equal 10--13M budgets, success rose from 1.23\% in L6 to 8.81\% in L7.
Near-ground tilt fell from \SI{5.89}{\degree} to \SI{2.98}{\degree}, angular rate from
\SI{25.2}{\degree\per\second} to \SI{9}{\degree\per\second}, and vertical contact speed from
\SI{8.88}{m/s} to \SI{7.50}{m/s}. Yet only 34\% of sampled trajectories moved toward the pad
centre. Step telemetry finally identified the cause: positive ``readiness'' was recorded on
about 78\% of steps because descending increased the altitude gate even when lateral error grew.

\subsection{L8--L9: separate learnability from realism}

L8 removed the altitude derivative from the potential and temporarily shortened engine startup,
minimum-run and cooldown delays. This was a deliberate curriculum over actuator dynamics: the
goal was first to make navigation, braking and four-foot contact discoverable. At the same
10--13M budget, L8 achieved 82.4\% success versus L7's 8.8\%. Recent successful L8 landings
reduced an initial \SI{20.1}{m} pad distance to \SI{1.41}{m}, contacted at \SI{2.26}{m/s} and
\SI{1.13}{\degree} tilt, and shut down propulsion correctly after impact.

The easier actuator also exposed a new strategy: nearly 30 ignitions per success. L9 restored
part of the timing constraint while initializing from L8. At matched 55--65\% difficulty, L9
slightly improved non-replay success and reduced hard touchdowns and restarts, but used more
fuel and missed the centre more often. Above \SI{50}{m}, more than 77\% of powered samples still
used all three engines. Thus timing continuation reduced cycling without producing the desired
single-engine navigation mode.

\subsection{L10--L11: reward the mission, then test its boundary}

L10 deliberately did not prescribe which engine to use. It started from fresh weights and paid
an efficiency term only after successful touchdown. Fuel, restart-equivalent fuel and additional
first ignitions affected this outcome term, while failures could not improve return by shutting
down early. Training was extended from 30M to 50M and finally 79.63M steps as the curriculum
continued to advance. Chapter~\ref{ch:results} shows that the same policy improved its landing
quality under increasingly difficult starts, although full-difficulty evaluation remained
outside the learned distribution.

L11 tested that hypothesis in a fresh 100M-step run. Minimum throttle was lowered from 57\% to
49.9\%, while startup delay, shutdown transient, minimum run and restart cooldown were set to
\SI{0.45}{s}, \SI{0.20}{s}, \SI{1.5}{s} and \SI{2.0}{s}. The lower minimum thrust made a longer
single-engine descent feasible; the timing constraints and larger success-only restart
equivalent made rapid cycling unattractive without rewarding failures for switching off.

The curriculum reached 100\% and stayed there. In the frozen 250-episode evaluation, L11
completed 212 stable landings, including 18 of 50 at exact full difficulty. Mean engine restarts
fell from L10's 18.46 to 1.56 per episode and mean first-contact vertical speed fell from
\SI{3.02}{m/s} to \SI{1.33}{m/s}. Extending the same run from 90M to 100M steps produced a
smaller but measurable final improvement: overall success rose from 82.8\% to 84.8\% and
full-difficulty success from 30\% to 36\%. The final policy therefore resolved most of the
cycling problem and produced the first successes at the previously unsolved boundary, while
leaving a narrower
vertical-braking and deployment-consistency problem for future work.

\section{Challenge 2: aerodynamics and physics}

The landing campaign was not only a reward-design exercise. Some policies exposed mistakes in
the world they were learning from. I had no previous experience building a rocket-physics system,
and the most dangerous errors were not spectacular failures but plausible-looking forces with
the wrong meaning. Axial drag worked during descent but pointed incorrectly during upward motion;
the first grid-fin model gave too much authority because it treated an open lattice too much like
a solid surface.

Watching training revealed the symptoms, but component tests made them reproducible. That led to
the layered model in Chapter~\ref{ch:physics}: force components are separated, dimensional
quantities and coefficients are exposed, and tests check direction and bounds where practical.
Where high-fidelity data was unavailable, I chose an explicit approximation rather than hidden
precision. The simulator remains an experimental learning environment, not a flight-qualified
predictor.

\section{Challenge 3: system integration}

As the physical and learning models grew, connecting them became at least as difficult as
implementing them. The recurring symptoms were unclear ownership, duplicated settings, startup
ordering problems and broken backward compatibility. A new actuator could change the policy
schema; a task could change spawning and termination; a UI edit could otherwise alter a value
after a model had already been selected.

Each of those failures produced one of the contracts in Chapter~\ref{ch:architecture}. Editable
drafts become immutable snapshots. Policy schemas are derived and compared. Resumes create
append-only revisions. Telemetry carries configuration identity. These mechanisms are not
decorative software engineering around the RL experiment; they are what make it possible to
tell whether two runs actually differ in the intended way.

\section{Remaining technical debt}

The resulting simulator is large for a single-author bachelor project, and the redesign is not
complete in every corner. Some partial classes remain long, older serialized assets require
migration care, and not every UI path has an automated test. The final public commit also remains
to be frozen. These limitations follow directly from the same story: every new physical or
learning capability created a software contract that the original prototype had never needed.
